El la realización de esta tarea se utilizó un solo computador de las siguientes especificaciones:

\begin{itemize}
    \item \textbf{Procesador:} AMD Ryzen 5 4500U con Radeon Graphics (2.38 GHz)
    \item \textbf{Memoria RAM:} 16 GB
    \item \textbf{Sistema Operativo:} Windows 11
    \item \textbf{Entorno de ejecución:} Subsistema de Windows para Linux (WSL) versión 2.6.3.0
    \item \textbf{Compilador:} GCC (g++)
\end{itemize}

Con respecto los datasets utilizados corresponden a los descritos en el enunciado de la tarea, pero es importante mencionar que para los algoritmos de ordenamiento no se ejecutó la totalidad de los casos más grandes, debido a que superaban las 4 horas de ejecución en  \texttt{ Quick Sort}, debido a que la tendencia de las curvas de rendimiento entorno a la memoria y tiempo ya eran claramente reconocibles con las dimensiones anteriores.

Finalmente, el tiempo de ejecución se cronometró con la librería \texttt{<chrono>}. Para evaluar la complejidad espacial y evadir la contaminación por el reciclaje de memoria de WSL, se aplicó una estimación algorítmica unificada en ambos experimentos: se determinó la memoria total sumando el peso de la estructura base (\texttt{std::vector}) y el espacio auxiliar dinámico demandado por cada algoritmo.